\subsection{Ring of types of quadratic forms.}

We shall assume, in this section, that A is a field of characteristic ≠ 2.

Given two quadratic forms Q, Q' on vector spaces V, V' over A, we shall call the tensor product of Q and Q', and denote by Q ⊗ Q', the quadratic form on V ⊗ V' whose associated bilinear form is the tensor product (§ 1, no. 9, def. 11) of the bilinear forms associated with Q and Q'. One sees readily that Q ⊗ Q' satisfies the relation

\[
(Q \otimes Q ^ {\prime}) (x \otimes x ^ {\prime}) = Q (x) Q ^ {\prime} (x ^ {\prime}) \quad (x \in V, x ^ {\prime} \in V ^ {\prime}).\tag{5}
\]

If Q and Q' are nondegenerate and of finite dimension, then so is Q ⊗ Q' (§ 1, no. 9, prop. 9).

Let Q, Q′, Q″ be quadratic forms on the vector spaces V, V′, V″. Using the canonical isomorphism from V ⊗ V' onto V' ⊗ V (resp. from (V ⊗ V′) ⊗ V″ onto V ⊗ (V′ ⊗ V″), from (V × V′) ⊗ V″ onto (V ⊗ V″) × (V′ ⊗ V″)), we see immediately that Q ⊗ Q' is equivalent to Q' ⊗ Q (resp. (Q ⊗ Q′) ⊗ Q″ to Q ⊗ (Q′ ⊗ Q″), (Q τ Q′) ⊗ Q″ to (Q ⊗ Q″) τ (Q′ ⊗ Q″)).

Let Q and Q' be two non-degenerate quadratic forms of finite dimensions. If Q is neutral, then so is Q⊗Q'. Indeed, let V, V' be the spaces of definition of Q, Q', with dimensions 2n and n', and let W be a totally singular subspace of V of dimension n (§ 4, no. 2); then W⊗V' is a totally singular subspace, and its dimension is half that of V⊗V'; from this one deduces, as in prop. 3, that Q⊗Q' is neutral. Likewise, Q⊗Q' is neutral whenever Q' is neutral.

From this we deduce that, if Q, Q', Q₁, Q₁' are nondegenerate quadratic forms of finite dimension over A, and if one assumes that θ(Q₁) = θ(Q) and θ(Q₁') = θ(Q'), then one has θ(Q₁ ⊗ Q₁') = θ(Q ⊗ Q'). Indeed, it suffices to verify this in the case where \(Q_{1} = Q \tau N\) and \(Q_{1}' = Q' \tau N'\) , N and \(N'\) being neutral forms; in this case \(Q_{1} \otimes Q_{1}'\) is equivalent to

\[
(Q \otimes Q ^ {\prime}) \tau (Q \otimes N ^ {\prime} \tau Q ^ {\prime} \otimes N \tau N \otimes N ^ {\prime})
\]

and the second parenthesis denotes a neutral form; this proves our assertion.

Now let \(\mathfrak{W}\) be the group of types of quadratic forms over A. Let us define, on the set \(\mathfrak{W}\) a second law of composition, denoted multiplicatively, by the formula

\[
\mathrm{TT} ^ {\prime} = \theta (\mathrm{T} \otimes \mathrm{T} ^ {\prime})\tag{6}
\]

\[
(\mathbf {T}, \mathbf {T} ^ {\prime} \text {   in   } \mathfrak {W}).
\]

It follows immediately from what we have just seen that this composition law is commutative, associative, and distributive with respect to addition. It admits an identity element, namely the type \(T_{1}\) of the quadratic form \(Q_{1}\) defined on the vector space A and such that \(Q_{1}(1)=1\) : indeed, by (5), \(Q_{1}\otimes Q=Q\) for every quadratic form Q. The additive group W, equipped with the multiplication just defined, is therefore a commutative ring with identity element; it is called the ring of types of quadratic forms of A (or the Witt ring of A, when no confusion is to be feared).

Remarks. --- 1) It is clear that, if a is an element of A*, then

\[
a. (\mathrm{TT} ^ {\prime}) = (a. \mathrm{T}) \mathrm{T} ^ {\prime} = \mathrm{T} (a. \mathrm{T} ^ {\prime})\tag{7}
\]

for any elements T, T' of W. One will note moreover that one has a.T = TaT, denoting by Ta the type of the quadratic form aQ1 on A.

\begin{enumerate}
\def\labelenumi{\arabic{enumi})}
\setcounter{enumi}{1}
\tightlist
\item
  Since A has characteristic \(\neq 2\) , every element T of \(\mathfrak{W}\) takes the form \(\sum_{i=1}^{n} a_i \cdot T_1\) where \(a_i \in A^*\) (no. 2, prop. 4). We have
\end{enumerate}

\[
(\sum_ {i = 1} ^ {n} a _ {i}. \mathrm{T} _ {1}) (\sum_ {j = 1} ^ {q} b _ {j}. \mathrm{T} _ {1}) = \sum_ {i, j} a _ {i} b _ {j}. \mathrm{T} _ {1} \quad (a _ {i}, b _ {j} \text {   in   } \mathrm{A} ^ {*}).\tag{8}
\]

\begin{enumerate}
\def\labelenumi{\arabic{enumi})}
\setcounter{enumi}{2}
\tightlist
\item
  Suppose that A is a maximal ordered field. Then the ring W is isomorphic to Z (prop. 5), the integer corresponding to the type of a form of signature \((s, t)\) being s - t (ibid.). Since the tensor product of two forms Q, \(Q'\) of signatures \((s, t)\) , \((s', t')\) is a form of dimension \((s + t)(s' + t')\) , it follows, by an elementary calculation, that the signature of \(Q \otimes Q'\) is \((ss' + tt', st' + ts')\) .
\end{enumerate}

